\mysection{Conclusion}
- tutorials not complete, different \\
- which model is better \\
- prospects, improvements, Recommendations \\
\\
\\ \\
Beispiele fürs referenzieren: \\
In \figref{FH-Logo} ist das HS München Logo zu sehen.

\begin{figure}[htbp]
\includegraphics[width=0.3\textwidth]{includes/MUASlogo}
\caption{FH-Logo}
\label{fig:FH-Logo}
\end{figure}

Oder auch eines Codes wie in  \listref{python_Code}.\\
\begin{lstlisting}[caption=Some python code, label=list:python_Code]

	bottleneck_path_2_bottleneck_values = {}


	def create_bottleneck_file(bottleneck_path, image_lists, label_name, index,
	                           image_dir, category, sess, jpeg_data_tensor,
	                           decoded_image_tensor, resized_input_tensor,
	                           bottleneck_tensor):
	  """Create a single bottleneck file."""
	  tf.logging.info('Creating bottleneck at ' + bottleneck_path)
	  image_path = get_image_path(image_lists, label_name, index,
	                              image_dir, category)
	  if not gfile.Exists(image_path):
	    tf.logging.fatal('File does not exist %s', image_path)
	  image_data = gfile.FastGFile(image_path, 'rb').read()
	  try:
	    bottleneck_values = run_bottleneck_on_image(
	        sess, image_data, jpeg_data_tensor, decoded_image_tensor,
	        resized_input_tensor, bottleneck_tensor)
	  except Exception as e:
	    raise RuntimeError('Error during processing file %s (%s)' % (image_path,
	                                                                 str(e)))
	  bottleneck_string = ','.join(str(x) for x in bottleneck_values)
	  with open(bottleneck_path, 'w') as bottleneck_file:
	    bottleneck_file.write(bottleneck_string)
\end{lstlisting}

Sectionrefs: In \secref{Methodological fundamentals} ist vieles noch nicht fertig. \\
SubSectionrefs: In \subsecref{Common Frameworks for Deep Learning Applications} wird dann näher auf den Inhalt eingegangen.\\
SubSubSectionrefs: In \subsubsecref{Tensorflow based on Python} gehts ans eingemachte.\\

Beispiele fürs zitieren:

Für einen noch besseren Überblick, kann das Buch von \citet{Butler2017} hinzugezogen werden.\\
Wenn in klammern und Seitenzahl \citep[p. 3]{Butler2017} \\

als compared, aber ohne Seitenzahl \citep[cmp.][]{Butler2017} \\
als compared mit Seitenzahl, das nd heißt "no date", da keine Jahrezahl vorhanden \citep[cmp.][p. 5]{Wang} \\

%Beispiele
%\citet{goossens93} -> Goossens et al. (1993)
%\citep{goossens93} -> (Goossens et al., 1993)
%
%\citet*{goossens93} -> Goossens, Mittlebach, and Samarin (1993)
%\citep*{goossens93} -> (Goossens, Mittlebach, and Samarin, 1993)